\documentclass[11pt]{article}

\usepackage[a4paper,margin=0.78in]{geometry}
\usepackage{graphicx}
\usepackage{booktabs}
\usepackage{array}
\usepackage{float}
\usepackage[hidelinks]{hyperref}
\usepackage{xcolor}
\usepackage{caption}
\usepackage{subcaption}
\usepackage{enumitem}
\usepackage{amsmath}

\graphicspath{{figures/}}
\setlength{\parindent}{0pt}
\setlength{\parskip}{0.55em}
\setlist[itemize]{topsep=0.25em,itemsep=0.15em,leftmargin=1.3em}
\setlist[enumerate]{topsep=0.25em,itemsep=0.15em,leftmargin=1.4em}

\title{Catch-up Report: Input Recovery from CW305 Power Traces}
\author{Khalifa University lab work notes}
\date{30 August 2026}

\begin{document}
\maketitle

\section*{Short update}

This note is for catching up with the professor. I was trying to recreate the
main idea of the Power2Picture and active power-template papers on our own
bench. The bench is not the same as the paper. I used the CW305 FPGA board and
ChipWhisperer-Lite. The paper used a different FPGA setup and stronger power
measurement conditions. So I adapted the victim design and the capture method.

The important update is that the earlier blocker was fixed. Before, after
programming, the board still looked like the stock AES design. This was not safe
for any final claim. After checking that M0, M1, and M2 were all high, and after
pressing USB RST/SW3, I reran the hardware checks. The custom design now passed
the check:

\begin{itemize}
  \item register pages 2, 3, and 4 read back as \texttt{0x00}, \texttt{0x00},
  and \texttt{0x00}, so it is not the stock AES readback;
  \item the hardware functional check passed all 676 convolution outputs against
  the CPU golden model;
  \item the new trace sets were captured from live CW305 and ChipWhisperer
  hardware, not from simulation;
  \item I added more tests after that: one active-template run with trained
  kernels, and two P2P feature/loss ablations;
  \item the active-template and Power2Picture-style runs reconstruct recognizable
  MNIST digits from the side-channel trace.
\end{itemize}

So the conclusion changed. I would no longer say the work is blocked by
programming mode. I would say the old issue was found and fixed, and now we have
a checked hardware rerun that can be shown as the main result.

\begin{figure}[H]
  \centering
  \includegraphics[width=\linewidth]{fig01_status_summary.png}
  \caption{The current outcome after the post-reset rerun. Normal result figures
  use color again. Only the flow-chart style diagrams are kept in black, white,
  and gray.}
  \label{fig:summary}
\end{figure}

\section*{What changed after pressing reset}

The older report had one serious warning. The board was giving this readback
after programming:

\begin{itemize}
  \item page 2 = \texttt{0x02},
  \item page 3 = \texttt{0x05},
  \item page 4 = \texttt{0x2e}.
\end{itemize}

These values match the stock CW305 AES core. If I continued from that state,
then the PC scripts could still run, but the FPGA would not be running my neural
network leakage design. This is dangerous because the software can make numbers
that look real, while the real board is doing another workload.

After the switch check and USB reset, I reran:

\begin{itemize}
  \item \path{host/cw305_leakage_func_check.py}
  \item \path{host/cw305_leakage_capture.py}
  \item \path{host/cw305_p2p_capture.py}
\end{itemize}

The new functional check printed:

\begin{quote}
\texttt{design\_signature\_pages\_2\_3\_4=[0, 0, 0]}\\
\texttt{PASS: 676 outputs bit-exact}
\end{quote}

This is the hardware check I now trust. It proves that the custom bitstream is
active enough for this experiment, and it proves that the FPGA output matches
the software golden model for the 26 by 26 first-layer window stream.

\begin{figure}[H]
  \centering
  \includegraphics[width=\linewidth]{fig03_register_identity_gate.png}
  \caption{The hardware identity check after the reset. The run is accepted only
  because the board does not answer like stock AES and the 676-output check
  passes.}
  \label{fig:identity}
\end{figure}

\section*{Evidence that the new data is real hardware data}

I captured two fresh trace sets.

The one-hot active-template set is:

\begin{itemize}
  \item folder: \path{host/traces_rerun_20260830_onehot_80},
  \item 80 MNIST images,
  \item 9 one-hot probe kernels per image,
  \item 5 repeated captures per kernel,
  \item 22,144 ADC samples per trace,
  \item live CW305/ChipWhisperer capture mode,
  \item ADC locked at the beginning and end,
  \item bitstream SHA-256 starts with
  \texttt{a5bb029e693174a3}.
\end{itemize}

The trained-kernel active-template set is:

\begin{itemize}
  \item folder: \path{host/traces_rerun_20260830_trained_80},
  \item 80 MNIST images,
  \item 9 trained convolution kernels per image,
  \item 5 repeated captures per kernel,
  \item 22,144 ADC samples per trace,
  \item live CW305/ChipWhisperer capture mode,
  \item the hardware functional check passed with zero mismatches.
\end{itemize}

The Power2Picture-style set is:

\begin{itemize}
  \item folder: \path{host/traces_rerun_20260830_p2p_2000},
  \item 2,000 MNIST images,
  \item one trained convolution kernel,
  \item one capture per image,
  \item 4 trace shards of 500 images,
  \item about 40 images per second during capture,
  \item live CW305/ChipWhisperer capture mode,
  \item the same bit-exact functional check before capture.
\end{itemize}

The validator passed both folders. It checked the manifest, live capture mode,
ADC lock, serial metadata, finite trace arrays, and the embedded functional
check. This matters because it separates the new run from the older pilot runs.

\begin{figure}[H]
  \centering
  \includegraphics[width=\linewidth]{fig16_rerun_trace_evidence.png}
  \caption{Fresh captured power signal data. The upper part shows the
  active-template trace segmentation and one power feature map. The lower part
  shows a trace from the P2P shard. The digit image is used after the attack for
  scoring and explanation.}
  \label{fig:fresh-trace}
\end{figure}

\section*{How the reconstruction uses the side-channel trace}

At attack time, I do not give the target image to the reconstruction algorithm.
The truth image is loaded only after the output exists, so I can score the
result. This point is important.

For the active-template method, the attack uses:

\begin{itemize}
  \item the captured power traces,
  \item timing knowledge: 676 valid 3 by 3 windows for one MNIST image,
  \item the profile template made from known profiling images,
  \item the selected one-hot probe kernels.
\end{itemize}

It does not use the attacked image or its label. In the clean one-hot rerun, I
profiled images \texttt{0700--0739} and attacked the held-out images
\texttt{0740--0779}. In the trained-kernel rerun, I profiled images
\texttt{0780--0819} and attacked \texttt{0820--0859}. I caught one bug here:
using \texttt{--start-index 40}
does not mean skip 40 captured files. The script interprets it as the MNIST file
index. That first precheck overlapped the profile set, so I did not use it in
the final table. I reran with \texttt{--start-index 740}, which gave the real
held-out set.

For the Power2Picture-style method, the model is trained with paired traces and
known images. This is supervised training, same general style as the paper. At
test time, the model receives only the trace features and outputs the image.
The test label and truth image are used only for scoring.

\begin{figure}[H]
  \centering
  \includegraphics[width=\linewidth]{fig05_algorithm_flow.png}
  \caption{Active/power-template reconstruction flow. The attack side uses
  traces and a profile. The truth image is outside the attack path.}
  \label{fig:algorithm}
\end{figure}

\section*{Fresh hardware results}

The fresh runs are the main results I would show first. They are smaller than a
final journal-scale experiment, but they are real checked hardware runs.

\begin{table}[H]
\centering
\caption{Fresh post-reset reconstruction results from live CW305 traces.}
\label{tab:fresh-results}
\small
\begin{tabular}{>{\raggedright\arraybackslash}p{0.28\linewidth}cccccc}
\toprule
Method & Test N & Pixel acc. & FG F1 & MSSIM & Rec. bin & Rec. gray \\
\midrule
Active template, one-hot held-out & 40 & 0.9455 & 0.7495 & 0.726 & 0.9750 & -- \\
Active template, trained kernels & 40 & 0.9600 & 0.7924 & 0.769 & 0.9000 & -- \\
P2P clock-abs, grad loss & 300 & 0.9503 & 0.7951 & 0.522 & 0.8933 & 0.9300 \\
P2P window-abs, grad loss & 300 & 0.9486 & 0.7796 & 0.508 & 0.8333 & 0.8533 \\
P2P clock-abs, MSE loss & 300 & 0.9537 & 0.8068 & 0.561 & 0.9033 & 0.9133 \\
\bottomrule
\end{tabular}
\end{table}

For the one-hot active-template result, the all-zero baseline bit accuracy was
0.8791. The recovered result reached 0.9455, so it is not only predicting black
background. The recognizer classified 39 out of 40 recovered images correctly.

The trained-kernel active-template run is closer to a normal first convolution.
It gave better pixel and structure scores than one-hot, but lower recognition:
36 out of 40 recovered images were classified correctly. This tells me that the
trained-kernel traces recover the stroke shape well, but some recovered digits
are shifted or broken in a way that affects the classifier.

For the P2P-style tests, all three models trained on the same 2,000 hardware
traces. The split was 1,500 train, 200 validation, and 300 test traces. The
clock-level feature kept more information than the window-level feature. The MSE
loss gave the best binary test result in this small run. The gradient loss gave
better gray-output recognition in the first clock-level run. So I would keep
both losses in the next large capture, and choose after a bigger test set.

\begin{figure}[H]
  \centering
  \includegraphics[width=\linewidth]{fig06_pilot_metrics.png}
  \caption{Fresh checked hardware results and P2P ablations. These entries come
  from the post-reset live CW305 trace folders.}
  \label{fig:metrics}
\end{figure}

\begin{figure}[H]
  \centering
  \includegraphics[width=0.70\linewidth]{fig15_rerun_onehot_heldout_examples.png}
  \caption{Fresh held-out active-template reconstruction examples. The attacked
  files did not contain the truth image or label.}
  \label{fig:active-examples}
\end{figure}

\begin{figure}[H]
  \centering
  \includegraphics[width=0.70\linewidth]{fig18_rerun_trained_heldout_examples.png}
  \caption{Fresh active-template examples using trained convolution kernels.
  This is a harder and more realistic setting than one-hot probes.}
  \label{fig:trained-rerun-examples}
\end{figure}

\begin{figure}[H]
  \centering
  \includegraphics[width=0.82\linewidth]{fig17_rerun_p2p_examples.png}
  \caption{Fresh P2P-style generator examples from live CW305 traces. The middle
  column is the direct gray generator output. The right column is the thresholded
  binary image.}
  \label{fig:p2p-rerun-examples}
\end{figure}

\begin{figure}[H]
  \centering
  \includegraphics[width=0.82\linewidth]{fig19_rerun_p2p_mse_examples.png}
  \caption{Fresh P2P examples from the clock-absolute feature with MSE loss.
  This was the best binary P2P result in the small ablation.}
  \label{fig:p2p-mse-examples}
\end{figure}

\section*{Older pilot results}

The older pilot results are still useful because they show the direction. But I
would not mix them with the fresh checked results without labelling them. Their
trace files pass file validation, but their manifests are not as strong as the
new ones. Some of them do not store the serials and embedded functional-check
object.

\begin{table}[H]
\centering
\caption{Older pilot outcomes. I keep these as background evidence, not as the
main final claim.}
\label{tab:old-results}
\begin{tabular}{>{\raggedright\arraybackslash}p{0.33\linewidth}ccccc}
\toprule
Method & N & Pixel acc. & FG F1 & MSSIM & Recognition \\
\midrule
Active template, trained kernels & 30 & 0.9676 & 0.8509 & 0.824 & 0.9333 \\
Active template, one-hot, Wei-scale & 200 & 0.9431 & 0.7779 & 0.704 & 0.9450 \\
Power2Picture-style generator & 500 & 0.9721 & 0.8909 & 0.803 & 0.9460 \\
Background detection, mean of 9 kernels & 30 & 0.8754 & 0.4859 & 0.382 & 0.8667 \\
\bottomrule
\end{tabular}
\end{table}

\begin{figure}[H]
  \centering
  \includegraphics[width=0.92\linewidth]{fig09_template_trained_examples.png}
  \caption{Older pilot active-template examples with trained convolution
  kernels. These are useful but not as strongly checked as the fresh rerun.}
  \label{fig:trained-examples}
\end{figure}

\begin{figure}[H]
  \centering
  \includegraphics[width=0.92\linewidth]{fig10_p2p_generator_examples.png}
  \caption{Older pilot Power2Picture-style generator examples. The generator can
  map trace features back to recognizable MNIST-like shapes.}
  \label{fig:p2p-old-examples}
\end{figure}

\begin{figure}[H]
  \centering
  \includegraphics[width=0.58\linewidth]{fig12_background_detection_examples.png}
  \caption{Older background-detection examples. This method was weaker because
  the adapted leakage source does not behave like a simple background detector.}
  \label{fig:bg}
\end{figure}

\section*{Why the previous implementation stalled and had low detection}

I think the old difficulty came from several issues together.

First, the hardware identity problem was real. If the board answers like stock
AES, then the PC can still capture traces and run algorithms, but the trace is
not from the neural-network leakage design. This was the biggest risk. The new
identity check fixes this by checking the register pages and the 676 hardware
outputs before capture.

Second, the first version tried to follow the original paper too directly. The
paper had better analog measurement conditions. Our CW-Lite setup has less
analog detail per clock cycle. Natural first-layer leakage can be hidden by
clock tree activity, static current, USB noise, and board-level noise. That is
why I used an adapted leakage core with a retained XNOR leak bank. It is still a
side-channel demo, but it must be described as an adapted controlled experiment,
not a direct clone of the old accelerator.

Third, the passive background method had weak leakage polarity for this design.
In the older paper, background pixels can appear as lower activity in the line
buffer or convolution datapath. In my adapted core, the power is closer to match
count between a window and a kernel. An all-background patch is not always the
lowest-power patch. So the simple background threshold stayed near the all-zero
baseline.

Fourth, I found a split issue during the rerun. The active-template script uses
\texttt{--start-index} as the image file index, not as the count offset inside
the capture folder. If I used \texttt{40}, the attack overlapped the first
profile images. I corrected this by using \texttt{740} for the held-out images.
This is small, but it is important for honest evaluation.

Fifth, one recognition subcommand still expected TensorFlow inside the
ChipWhisperer Python environment. That environment does not have TensorFlow.
The compact scorer used the NumPy recognizer, so the reported recognition
numbers are still valid. But this is another reason why the scripts should keep
the TensorFlow-free scoring path as the default.

\section*{How this compares with the original papers}

The research question is the same: can power traces from the first layer reveal
the private input image?

From this bench, the answer is yes, in a controlled and adapted setup. The fresh
run recovered recognizable digit images from real FPGA power traces. The
active-template method reached 97.5\% recovered recognition on 40 held-out
images. The P2P-style method reached 89.33\% binary recognition and 93.0\%
gray-output recognition on 300 test traces.

But the setup is not identical to the paper:

\begin{itemize}
  \item the paper target and our FPGA target are different;
  \item the paper measurement setup is stronger than CW-Lite;
  \item my active-template run uses one-hot probe kernels to make the profile
  easier to learn;
  \item my victim core includes an adapted leakage bank to make the leakage
  visible on this bench;
  \item the P2P rerun used only 2,000 traces, not a very large dataset.
\end{itemize}

\begin{figure}[H]
  \centering
  \includegraphics[width=\linewidth]{fig07_reference_vs_our_setup.png}
  \caption{The work keeps the same research question as the papers, but the
  hardware and leakage source are adapted for CW305 and ChipWhisperer-Lite.}
  \label{fig:compare}
\end{figure}

So I would present this as a strong hardware reproduction in spirit, not as a
bit-for-bit reproduction. The main research value for our work is the practical
adaptation:
showing input recovery on the available CW305/ChipWhisperer bench, with a strict
hardware identity check, trace-only evaluation bundles, and both template-based
and generator-based recovery.

\section*{What I would do next}

If we continue this toward a stronger paper, I would do these next steps:

\begin{enumerate}
  \item Add a build-ID register in the FPGA design, so every trace folder can
  prove the exact design name, version, dwell cycles, and leak-lane count.
  \item Repeat the active-template run with more held-out images and trained
  kernels, not only one-hot kernels.
  \item Scale the P2P capture to 10,000 or more traces and train longer. The
  gray generator output already looks useful, so more data should help.
  \item Add ablation tests: no leak bank, weaker leak bank, different dwell
  cycles, different gain, and different clock frequencies.
  \item Test stability across time and temperature, because the P2P generator
  may learn capture drift if the split is not random.
  \item Keep the same rule: no result is accepted unless live capture manifest,
  not-AES identity check, and 676-output functional check all pass.
\end{enumerate}

\begin{figure}[H]
  \centering
  \includegraphics[width=\linewidth]{fig08_next_hardware_checklist.png}
  \caption{The hardware checklist that should stay in the run order for future
  runs. It prevents the PC from training on traces from the wrong FPGA design.}
  \label{fig:next}
\end{figure}

\section*{My final view}

My view is that the experiment now has a good result. It is not only
simulation and it is not only host-PC reconstruction. The fresh runs used live
CW305 traces captured with ChipWhisperer. The board passed the custom-design
check after the switch/reset fix. Then the reconstruction methods recovered
recognizable digit images from the traces.

I would tell the professor:

\begin{quote}
I found why the previous run was unsafe to claim: the board could still answer
as stock AES after programming. After checking the mode switches and pressing
USB reset, the custom design check passed. I reran the capture on real CW305
hardware. The one-hot active-template method reached 97.5\% recognition on 40
held-out images. The trained-kernel template reached 90.0\%. For the P2P-style
generator, the best binary result was 90.33\%, and the best gray-output result
was 93.0\% on 300 test traces. The method is adapted for our ChipWhisperer
setup, but it keeps the same idea: first-layer power traces can leak enough
information to recover the input image.
\end{quote}

This is the point where the next student can continue. The next work should
scale the run, add the build-ID register, and compare the adapted leakage core
against a more natural accelerator implementation.

\section*{Files behind this report}

The shortened P2P summary names in this list are under
\path{attack/rerun_20260830/}.

\small
\begin{itemize}
  \item Active-template fresh trace folder:
  \path{host/traces_rerun_20260830_onehot_80}.
  \item Active-template fresh score:
  \path{attack/rerun_20260830/score_onehot_80_eval40_heldout.json}.
  \item Active-template detailed score:
  \path{score_summary.json} inside the one-hot held-out result folder.
  \item Trained-kernel fresh trace folder:
  \path{host/traces_rerun_20260830_trained_80}.
  \item Trained-kernel fresh score:
  \path{attack/rerun_20260830/score_trained_80_eval40_heldout.json}.
  \item P2P fresh trace folder:
  \path{host/traces_rerun_20260830_p2p_2000}.
  \item P2P clock-abs gradient-loss summary:
  \path{p2p_2000_clockabs/summary.json}.
  \item P2P window-abs gradient-loss summary:
  \path{p2p_2000_windowabs/summary.json}.
  \item P2P clock-abs MSE-loss summary:
  \path{p2p_2000_clockabs_mse/summary.json}.
  \item Active-template fresh validation output:
  \path{validation_rerun_onehot_80.json}.
  \item Trained-kernel fresh validation output:
  \path{validation_rerun_trained_80.json}.
  \item P2P fresh validation output:
  \path{validation_rerun_p2p_2000.json}.
  \item Real stock AES control trace metadata:
  \path{report/hardware_aes_control_trace_pass_20260829.json}.
  \item Older run summary:
  \path{HARDWARE_RUN_2026-08-27.md}.
  \item Older result table:
  \path{report/results_table_2026-08-27.md}.
\end{itemize}

\end{document}
